\documentclass[12pt]{article}
\usepackage{amsmath, amssymb, amsthm, bm}
\usepackage[margin=1in]{geometry}
\usepackage[skip=0.6em]{parskip}

\newtheorem{proposition}{Proposition}
\theoremstyle{remark}
\newtheorem{remark}{Remark}

\begin{document}

\begin{center}
{\Large Extension plan: redundant moments in a\\ studentized Gaussian limit experiment}\\[3pt]
{\small Ivan A. Canay :: August 27, 2026}
\end{center}

\medskip

\section*{1. Motivation and limit experiment}

Consider a moment inequality problem based on a known vector-valued function
$m(W,\theta)\in\mathbf R^d$.  The null hypothesis at a fixed value of $\theta$ is
\[
E[m(W,\theta)]\le0~.
\]
Under the usual local asymptotic approximation, the vector of studentized sample
moments converges to
\[
X\sim N_d(\mu,\Omega)~,
\]
where $\Omega$ is a positive-definite correlation matrix.  The correlation matrix is
part of the limit experiment and should be left unrestricted.  Restrictions such as
identity covariance or equicorrelation may be useful examples, but they do not capture
the limit experiments generated by the moment inequality models of interest.

Partition the moment indices into two nonempty sets.  Let $\mathcal C$ denote the
\emph{core} moments, which contain all identifying information, and let $\mathcal R$
denote the moments that are redundant for identification.  Write
\[
\mu=(\mu_{\mathcal C}',\mu_{\mathcal R}')'~.
\]
The proposed restriction on the standardized mean vector is
\begin{equation}\label{eq:redundancy}
\text{for every }r\in\mathcal R,\quad
\mu_r\le\mu_c\text{ for some }c\in\mathcal C~.
\end{equation}
Equivalently, define
\[
\Lambda
=
\left\{\mu\in\mathbf R^d:
\mu_r\le\max_{c\in\mathcal C}\mu_c
\text{ for every }r\in\mathcal R
\right\}~.
\]
The restricted null and alternative are
\[
\Lambda_0=\Lambda\cap(-\infty,0]^d~,
\qquad
\Lambda_1=\Lambda\setminus\Lambda_0
=
\left\{\mu\in\Lambda:
\max_{c\in\mathcal C}\mu_c>0
\right\}~.
\]
The last equality is the formal content of identification redundancy.  A redundant
moment cannot have a positive mean unless at least one core moment also has a positive
mean.  Therefore the core moments determine whether the null is violated, even though
the redundant moments may still contain useful sampling information.

\section*{2. Why the restriction is imposed after studentization}

Condition~\eqref{eq:redundancy} is a restriction on the standardized local means.  This
point should be stated explicitly.  If $\delta_j$ denotes the local drift of the
unstandardized $j$th moment and $\sigma_j$ is its standard deviation, then
\[
\mu_j=\frac{\delta_j}{\sigma_j}~.
\]
An ordering of the unstandardized drifts, such as $\delta_r\le\delta_c$, becomes
\[
\sigma_r\mu_r\le\sigma_c\mu_c~,
\]
and does not in general imply $\mu_r\le\mu_c$.  The standardized restriction is
therefore a substantive feature of the model, not a normalization that follows
automatically from redundancy of the original moments.  It is also the scale-invariant
formulation: multiplying a moment by a positive constant changes neither its
studentized statistic nor its standardized local mean.

An applied version of the paper should explain when the structure of the moment
functions implies~\eqref{eq:redundancy} after studentization.  This is the link between
the Gaussian experiment and the moment inequality models that motivate the analysis.

\section*{3. Feasible normals under arbitrary correlation}

Let $A\subseteq\mathbf R^d$ be the acceptance region of the nonrandomized test
\[
\phi_A(x)=I\{x\notin A\}~.
\]
Suppose that $A$ is closed, convex, and coordinatewise lower, and that
\[
P_0(X\in A)=1-\alpha~.
\]
For a closed convex set $C$, write
\[
h_C(a):=\sup_{x\in C}a'x
\]
for its support function.  Since $A$ is coordinatewise lower, $h_A(a)<\infty$ for a
nonzero direction $a$ only if $a\ge0$.

The natural parameter in the Gaussian experiment is $\vartheta=\Omega^{-1}\mu$.
Consequently, a supporting normal $a$ generates the natural-parameter ray
$\vartheta=\lambda a$ and the corresponding mean ray
\[
\mu=\lambda\Omega a~.
\]
This suggests the feasible-normal set
\begin{equation}\label{eq:feasible-set}
\mathcal F(\Omega,\Lambda)
:=
\{a\in\mathbf R^d:a\ge0,\ \Omega a\in\Lambda\}~.
\end{equation}

\begin{proposition}[Feasible-normal sufficiency]\label{prop:sufficiency}
Let $X\sim N_d(\mu,\Omega)$, where $\Omega$ is a positive-definite correlation matrix,
and let $\Lambda_0$ and $\Lambda_1$ be defined above.  Suppose that $A$ is closed,
convex, and coordinatewise lower, satisfies $P_0(X\in A)=1-\alpha$, and
\[
a\ne0,\quad h_A(a)<\infty
\quad\Longrightarrow\quad
a\in\mathcal F(\Omega,\Lambda)~.
\]
Then $\phi_A$ is $\alpha$-admissible against all level-$\alpha$ tests, including
randomized competitors.
\end{proposition}

The argument is the same Lehmann--Romano ray argument used in the bivariate note.  If
$a$ is a finite supporting normal, the mean ray $\lambda\Omega a$ belongs to $\Lambda$.
It cannot remain in the null: if $\Omega a\le0$, then
\[
a'\Omega a\le0~,
\]
which contradicts positive definiteness because $a\ne0$.  Thus $\Omega a$ has a
positive coordinate.  Condition~\eqref{eq:redundancy} then implies that some core
coordinate is positive, so $\lambda\Omega a\in\Lambda_1$ for every $\lambda>0$.
Finally, coordinatewise lowerness makes $\mu=0$ a least favorable null point and gives
exact size there.  These observations verify the ray and size conditions.

\section*{4. An admissible test that uses redundant moments}

Proposition~\ref{prop:sufficiency} gives a simple construction that works for every
positive-definite correlation matrix.  Fix any $c_0\in\mathcal C$, and let
$b\in\mathbf R^d$ satisfy
\[
b_c=0\quad(c\in\mathcal C),
\qquad
b_r>0\quad(r\in\mathcal R)~.
\]
For $\varepsilon>0$, define
\[
a_\varepsilon=e_{c_0}+\varepsilon b~.
\]
At $\varepsilon=0$,
\[
(\Omega e_{c_0})_{c_0}=1,
\qquad
(\Omega e_{c_0})_r=\Omega_{r c_0}<1
\quad\text{for every }r\in\mathcal R~.
\]
The strict inequality follows from positive definiteness and the fact that $\Omega$ is
a correlation matrix.  Since $\mathcal R$ is finite, continuity implies that, for all
sufficiently small $\varepsilon>0$,
\[
(\Omega a_\varepsilon)_r
<
(\Omega a_\varepsilon)_{c_0}
\quad\text{for every }r\in\mathcal R~.
\]
Hence $a_\varepsilon\in\mathcal F(\Omega,\Lambda)$.

Let $z_{1-\alpha}=\Phi^{-1}(1-\alpha)$ and consider
\begin{equation}\label{eq:augmented-test}
\phi_\varepsilon(x)
=
I\left\{
a_\varepsilon'x>
z_{1-\alpha}\sqrt{a_\varepsilon'\Omega a_\varepsilon}
\right\}~.
\end{equation}
This test assigns strictly positive weight to every redundant moment.  Its acceptance
region is a closed, convex, coordinatewise lower halfspace, and its only finite normal
is $a_\varepsilon$.  Proposition~\ref{prop:sufficiency} therefore gives the following
conclusion.

\begin{proposition}[Admissible use of redundant moments]\label{prop:existence}
For every positive-definite correlation matrix $\Omega$, there exists
$\varepsilon>0$ such that the test in~\eqref{eq:augmented-test} has exact size $\alpha$
and is $\alpha$-admissible against all level-$\alpha$ tests.  The test places strictly
positive weight on every moment in $\mathcal R$.
\end{proposition}

This result makes the paper's main point directly.  Moments that are redundant for
identification need not be irrelevant for inference.  They can enter a test that is
globally admissible in the Gaussian limit experiment, without imposing any special
structure on the correlation matrix.

\section*{5. Proposed structure of the argument}

The paper can make the distinction between identification and inference in three
steps.

First, the existing local-power calculations compare tests based only on the core
moments with tests that also use redundant moments.  The crossing power curves show
that dropping redundant moments is not uniformly desirable.

Second, Propositions~\ref{prop:sufficiency} and~\ref{prop:existence} provide the
decision-theoretic result in the general studentized Gaussian experiment.  The main
claim requires neither identity covariance nor equicorrelation: for every correlation
matrix, some globally admissible test uses the redundant moments.

Third, the bivariate identity-covariance characterization explains why the manner in
which redundant moments enter still matters.  The short and gated-sum tests use only
feasible normals and are globally admissible, whereas the long rectangle uses an
infeasible vertical normal and is globally inadmissible.  This sharp example prevents
the general result from being misread as saying that every augmented test is well
designed.

A full necessity characterization for arbitrary $\Omega$ is not needed for this
argument.  The next tasks are instead to verify Proposition~\ref{prop:sufficiency}
against the exact Lehmann--Romano statement, connect condition~\eqref{eq:redundancy}
to the motivating moment inequality models after studentization, and decide whether
the admissible construction can be embedded in a familiar class of moment inequality
statistics.

\begin{thebibliography}{9}

\bibitem{LehmannRomano2005}
Lehmann, E. L. and Romano, J. P. (2005).
\textit{Testing Statistical Hypotheses}. Third edition. Springer.

\end{thebibliography}

\end{document}
